

\para{Practical implications.} The compute-optimal allocation frontier (\autoref{fig:allocation_frontier}) provides concrete guidance for data-constrained practitioners. Given a fixed data budget, there is a compute threshold beyond which repeating data further is counterproductive. Past this threshold, compute is better allocated to model capacity. Combined with the finding that strong weight decay reduces the overfitting coefficient $P$ by approximately 70\%, practitioners in data-constrained settings have two complementary levers: increasing regularization strength and choosing the right model size--epoch tradeoff.

\para{Limitations.} Our study spans model sizes up to 1B parameters and repetition counts up to 16 epochs; we cannot confirm that the fitted exponents hold at the frontier scales used in practice. Our scaling law form may not extrapolate to extreme repetition counts beyond those covered in our study. Notably, our law doesn't encapsulate phenomena such as double descent \citep{Nakkiran2020Deep}. We test only two weight decay values ($\lambda \in \{0.1, 1.0\}$) and fit independent scaling sweeps to provide prescriptive recommendations. Extending our law to incorporate regularization strength directly is an interesting direction for future work.
